\chapter{Bitvector Slices\label{chap:BitvectorSlicing}}
\hypertarget{bitsliceterm}{}
Bitvector slices allow extracting a contiguous sub-range of bits out of a given bitvector to produce
another bitvector.
Lists of bitvector slices can extract several sub-ranges and arrange them in different orders in the
output bitvector.

\ExampleDef{Bitvector Slices}
\listingref{BitvectorSlices} shows different ways of expressing bitvector slices
and how they relate to one another in terms of the order of bits they represent.
\ASLListing{Examples of bitvector slices and operations on slices}{BitvectorSlices}{\definitiontests/Bitvector_slices.asl}

\ChapterOutline

%%%%%%%%%%%%%%%%%%%%%%%%%%%%%%%%%%%%%%%%%%%%%%%%%%%%%%%%%%%%%%%%%%%%%%%%%%%%%%%%
\FormalRelationsDef[defines the formal relations for bitvector slices;]{Bitvector Slices}
\paragraph{Syntax:} Bitvector slices are grammatically derived from $\Nslice$.

\paragraph{Abstract Syntax:} Bitvector slices are derived in the abstract syntax from $\slice$,
and generated by $\buildslice$ (see \ASTRuleRef{Slice}).

\paragraph{Typing:} Bitvector slices are annotated by $\annotateslice$ (see \TypingRuleRef{Slice}).

\paragraph{Semantics:} Bitvector slices are evaluated by $\evalslice$.

%%%%%%%%%%%%%%%%%%%%%%%%%%%%%%%%%%%%%%%%%%%%%%%%%%%%%%%%%%%%%%%%%%%%%%%%%%%%%%%%
\SyntaxDef[defines the syntax of bitvector slices;]{Bitvector Slices}
\begin{flalign*}
\Nslice \derives \ & \Nexpr &\\
            |\  & \Nexpr \parsesep \Tcolon \parsesep \Nexpr &\\
            |\  & \Nexpr \parsesep \Tpluscolon \parsesep \Nexpr &
\end{flalign*}

%%%%%%%%%%%%%%%%%%%%%%%%%%%%%%%%%%%%%%%%%%%%%%%%%%%%%%%%%%%%%%%%%%%%%%%%%%%%%%%%
\AbstractSyntaxDef[defines the abstract syntax of bitvector slices;]{Bitvector Slices}
\RenderTypes[remove_hypertargets]{untyped_slice}

\ASTRuleDef{Slice}
\NotImportedToASLSpecYet{
\hypertarget{build-slice}{}
The function
\[
  \buildslice(\overname{\parsenode{\Nslice}}{\vparsednode}) \;\aslto\; \overname{\slice}{\vastnode}
\]
transforms a parse node for a slice $\vparsednode$ into an AST node for a slice $\vastnode$.
} % END_OF_UNIMPORTED_RELATION


\begin{mathpar}
\inferrule[single]{}{
  \buildslice(\Nslice(\punnode{\Nexpr})) \astarrow
  \overname{\SliceSingle(\astof{\vexpr})}{\vastnode}
}
\end{mathpar}

\begin{mathpar}
\inferrule[range]{
  \buildexpr(\veone) \astarrow \astversion{\veone}\\
  \buildexpr(\vetwo) \astarrow \astversion{\vetwo}
}{
  \buildslice(\Nslice(\namednode{\veone}{\Nexpr}, \Tcolon, \namednode{\vetwo}{\Nexpr})) \astarrow
  \overname{\SliceRange(\astversion{\veone}, \astversion{\vetwo})}{\vastnode}
}
\end{mathpar}

\begin{mathpar}
\inferrule[length]{
  \buildexpr(\veone) \astarrow \astversion{\veone}\\
  \buildexpr(\vetwo) \astarrow \astversion{\vetwo}
}{
  \buildslice(\Nslice(\namednode{\veone}{\Nexpr}, \Tpluscolon, \namednode{\vetwo}{\Nexpr})) \astarrow
  \overname{\SliceLength(\astversion{\veone}, \astversion{\vetwo})}{\vastnode}
}
\end{mathpar}

%%%%%%%%%%%%%%%%%%%%%%%%%%%%%%%%%%%%%%%%%%%%%%%%%%%%%%%%%%%%%%%%%%%%%%%%%%%%%%%%
\TypeRulesDef[defines the type rules for bitvector slices;]{Bitvector Slices}
\TypingRuleDef{Slice}
\RenderRelation{annotate_slice}

\ExampleDef{Annotating Slices}
The slices in
\listingref{semantics-slicesingle},
\listingref{semantics-slicerange}, and
\listingref{semantics-slicelength}
are all well-typed.

\RenderProseAndFormally{annotate_slice}

\TypingRuleDef{SlicesWidth}
\RenderRelation{slices_width}

\ExampleDef{The Width of a List of Slices}
In \listingref{eval-slices}
The width of the list of slices \verb|[2, 7:5, 0+:3]| is $7$.

\RenderProseAndFormally{slices_width}

\TypingRuleDef{SliceWidth}
\RenderRelation{slice_width}

\ExampleDef{The Width of a Slice}
In \listingref{semantics-slicesingle},
the width of the slice \texttt{[2]} is $1$.

In \listingref{semantics-slicerange},
the width of the slice \texttt{4:2} is $3$.

In \listingref{semantics-slicelength},
the width of the slice \texttt{2+:3} is $3$.

\RenderProseAndFormally{slice_width}

\TypingRuleDef{SymbolicConstrainedInteger}
\RenderRelation{annotate_symbolic_constrained_integer}

\ExampleDef{Annotating Symbolically Evaluable Constrained Integer Expressions}
In \listingref{typing-annotatesymbolicallyevaluableexpr}, all of the symbolically evaluable
expressions (as noted by the comments) are not constrained as their \underlyingtypesterm{} are
the \unconstrainedintegertypeterm.

On the other hand, in \listingref{check-constrained-integer}
all of the \rhsexpressions{} of assignments are both symbolically
evaluable \underline{and} their \underlyingtypesterm{} are \constrainedintegerterm{} types,
since they consist of \staticallyevaluableterm{} expressions (literals) and
the parameter \verb|N|.

\RenderProseAndFormally{annotate_symbolic_constrained_integer}
\CodeSubsection{\SymbolicConstrainedIntegerBegin}{\SymbolicConstrainedIntegerEnd}{../Typing.ml}

%%%%%%%%%%%%%%%%%%%%%%%%%%%%%%%%%%%%%%%%%%%%%%%%%%%%%%%%%%%%%%%%%%%%%%%%%%%%%%%%
\SemanticsRulesDef[defines the dynamic semantics of bitvector slices;]{Bitvector Slices}
\SemanticsRuleDef{Slice}
\RenderRelation{eval_slice}

\ExampleDef{Evaluating Slices}
In \listingref{semantics-slicesingle},
the slice \texttt{[2]} evaluates to \texttt{(2, 1)}, that is, the slice of
length 1 starting at index 2.
\ASLListing{A single expression slice}{semantics-slicesingle}{\semanticstests/SemanticsRule.SliceSingle.asl}

In \listingref{semantics-slicerange},
the slice \texttt{4:2} evaluates to \texttt{(2, 3)}.
\ASLListing{A range slice}{semantics-slicerange}{\semanticstests/SemanticsRule.SliceRange.asl}

In \listingref{semantics-slicelength},
the slice \texttt{2+:3} evaluates to \texttt{(2, 3)}.
\ASLListing{A slice by length}{semantics-slicelength}{\semanticstests/SemanticsRule.SliceLength.asl}

\RenderProseAndFormally{eval_slice}

%%%%%%%%%%%%%%%%%%%%%%%%%%%%%%%%%%%%%%%%%%%%%%%%%%%%%%%%%%%%%%%%%%%%%%%%%%%%%%%%
\FormalRelationsDef[defines the formal relations for lists of bitvector slices;]{Lists of Bitvector Slices}
\paragraph{Syntax:} Lists of bitvector slices are grammatically derived from $\Nslices$.

\paragraph{Abstract Syntax:} Lists of bitvector slices are represented by a list of $\slice$
  AST nodes, and are generated by $\buildslices$.

\paragraph{Typing:} Lists of bitvector slices are annotated by $\annotateslices$\\ (see \TypingRuleRef{Slices}).

\paragraph{Semantics:} Lists of bitvector slices are evaluated by $\evalslices$\\ (see \SemanticsRuleRef{Slices}).

%%%%%%%%%%%%%%%%%%%%%%%%%%%%%%%%%%%%%%%%%%%%%%%%%%%%%%%%%%%%%%%%%%%%%%%%%%%%%%%%
\SyntaxDef[defines the syntax of lists of bitvector slices;]{Lists of Bitvector Slices}
\begin{flalign*}
\Nslices \derives \ & \Tlbracket \parsesep \ClistOne{\Nslice} \parsesep \Trbracket &
\end{flalign*}

%%%%%%%%%%%%%%%%%%%%%%%%%%%%%%%%%%%%%%%%%%%%%%%%%%%%%%%%%%%%%%%%%%%%%%%%%%%%%%%%
\AbstractSyntaxDef[defines the abstract syntax of lists of bitvector slices;]{Lists of Bitvector Slices}
\ASTRuleDef{Slices}
\NotImportedToASLSpecYet{
\hypertarget{build-slices}{}
The function
\[
  \buildslices(\overname{\parsenode{\Nslices}}{\vparsednode}) \;\aslto\; \overname{\KleenePlus{\slice}}{\vastnode}
\]
transforms a parse node for a list of slices $\vparsednode$ into an AST node for a list of slices $\vastnode$.
} % END_OF_UNIMPORTED_RELATION


\begin{mathpar}
\inferrule{
  \buildclist[\buildslice](\vslices) \astarrow \vsliceasts
}{
  \buildslices(\Nslices(\Tlbracket, \namednode{\vslices}{\ClistOne{\Nslice}}, \Trbracket)) \astarrow
  \overname{\vsliceasts}{\vastnode}
}
\end{mathpar}

%%%%%%%%%%%%%%%%%%%%%%%%%%%%%%%%%%%%%%%%%%%%%%%%%%%%%%%%%%%%%%%%%%%%%%%%%%%%%%%%
\TypeRulesDef[defines the type rules for lists of bitvector slices;]{Lists of Bitvector Slices}
\TypingRuleDef{Slices}
\RenderRelation{annotate_slices}

See \ExampleRef{Bitvector Slices} for examples of well-typed bitvector slices.

\ExampleDef{An Ill-typed Bitvector Slice}
In \listingref{ImpureBitvectorSlice}, the bit vector slice on \verb|x| is ill-typed, since the expression \\
\verb|side_effecting()| is not \pureterm{}.
\ASLListing{An ill-typed bitvector slice}{ImpureBitvectorSlice}{\typingtests/TypingRule.AnnotateSlices.bad-impure.asl}

\RenderProseAndFormally{annotate_slices}

%%%%%%%%%%%%%%%%%%%%%%%%%%%%%%%%%%%%%%%%%%%%%%%%%%%%%%%%%%%%%%%%%%%%%%%%%%%%%%%%
\SemanticsRulesDef[defines the dynamic semantics rules for lists of bitvector slices;]{Lists of Bitvector Slices}
\SemanticsRuleDef{Slices}
\RenderRelation{eval_slices}

\ExampleDef{Evaluating a List of Slices}
In \listingref{eval-slices}, evaluating the list of slices \verb|[2, 7:5, 0+:3]|
yields the list of ranges $[(2,1), (5,3), (0,3)]$.

\ASLListing{Evaluating a list of slices}{eval-slices}{\semanticstests/SemanticsRule.Slices.asl}

\RenderProseAndFormally{eval_slices}
\CodeSubsection{\EvalSlicesBegin}{\EvalSlicesEnd}{../Interpreter.ml}
